\section{Conclusion and Future Work}

This dissertation introduced a data-efficient pipeline for building new image quality metrics that scales from a small set of human-rated distortions to large unlabeled image sets, and that can either stop at a full-reference fusion metric or continue to a no-reference predictor, each optimized for the target application.

The pipeline begins with a compact but carefully curated subjective study. Human scores are screened and debiased to minimize observer and content effects. From there, a diverse set of full-reference measures is harmonized and fused into a single perceptual target. This stage yields \acs{pMOS-SSIM+}, a fused metric that aligns with human opinion more closely than any individual method and that can be adopted immediately wherever references are available.

Crucially, the same fused metric enables expansion beyond the limited human-rated subset. By assigning pseudo-labels to large collections without references, the pipeline supervises the learning of no-reference models that inherit the perceptual fidelity of the fusion while gaining real-world reach. Two outcomes illustrate the range of this second stage: \acs{pMOS-StegaFace}, specialized for steganographically affected facial images, and \acs{pMOS-NoisyTID}, designed to generalize across diverse content and distortion types. Both surpass established no-reference baselines, detect subtle degradations that conventional approaches often miss, and maintain strong agreement with human judgments. Throughout, fairness safeguards are integrated at the data and model levels so that quality estimates remain consistent across demographic groups.

The broader implication is a practical recipe for building metrics where exhaustive subjective annotation is not feasible. With only a small number of human-rated distortions to set direction, the pipeline first concentrates perceptual signal through fusion and then amplifies it across scale via pseudo-labels, producing metrics that can be deployed either as a high-fidelity full-reference standard or as no-reference models tuned to domain constraints.

Future work follows four paths. First, extend the seed set of human-rated distortions to cover additional subtle or task-critical mechanisms, and study how few new ratings are needed for the fusion to transfer reliably. Second, strengthen generalization of the no-reference stage with improved training curricula and data augmentation, targeting challenging capture conditions and authentic mixtures of artifacts. Third, integrate calibrated uncertainty and active sampling so that the system can request new human ratings only where they add the most value, further reducing annotation cost. Fourth, complete end-to-end deployment studies in biometric workflows, including threshold setting, early-warning triggers for tampering, and continuous monitoring for fairness across age, gender, and skin tone.

By unifying a compact subjective core, a robust fusion target, and scalable no-reference learning, this work delivers three concrete metrics: \acs{pMOS-SSIM+}, \acs{pMOS-StegaFace}, and \acs{pMOS-NoisyTID}, and, more importantly, a general pipeline to build the next ones. This foundation enables more reliable, fair, and security-aware image quality control in practice, raising the standard for how future systems screen and maintain image quality in alignment with human perception.

% We proposed a Full-to-No-Reference framework for \acs{FIQA} that predicts image quality in the absence of reference images by leveraging a pseudo-\acs{MOS} supervision strategy. Our method first trains a \acs{FR} fusion model to regress human perceptual judgments on a labeled subset (10\% of the dataset), generating pseudo-\acs{MOS} labels for a larger unlabeled dataset. These forged labels are then used to train a deep \acs{NR} regressor, enabling quality prediction from distorted images alone. This two-stage pipeline effectively bridges the gap between fully supervised \acs{FR}-\acs{IQA} and reference-free \acs{NR}-\acs{IQA} approaches.

% Beyond the development of our \acs{NR}-\acs{IQA} metric, the proposed framework offers a flexible foundation for constructing a variety of task-specific models. By enabling scalable, perceptually grounded supervision with limited ground-truth annotations, our approach can facilitate quality-aware training in applications such as \acs{GANs}, forensic imaging, and domain-adapted biometric pipelines.

